
\section{Conclusion}
\label{Conclusion}

We presented SECDA-DSE, an LLM-guided framework for FPGA accelerator design space exploration built upon the SECDA ecosystem, and extended its evaluation through end-to-end FPGA execution. SECDA-DSE combines structured DSE, retrieval-augmented reasoning, Chain-of-Thought prompting, and iterative hardware evaluation to generate accelerator designs from natural language specifications. 
Specifically, we generated and executed element-wise vector multiplication, 2D convolution, and matrix transpose accelerators targeting a Xilinx Zynq-7000 FPGA platform. The generated designs successfully progressed through the FPGA implementation flow and demonstrated workload-aware hardware behavior across different compute and memory access patterns. 
Our results provide further evidence that combining LLM-guided reasoning with structured FPGA design workflows can support adaptive accelerator generation across diverse workloads while reducing manual design effort and advancing more automated hardware design methodologies. The successful FPGA execution of all generated accelerators further demonstrates the ability of SECDA-DSE to adapt across different workload characteristics.
